%% kapa: exam=1 round=26 year=2015 period=1 form=A subject=economics subject_ko=경제원론 seq=30 number=70
%% source: Q-Net 감정평가사 기출문제 (cst00309) artlSeq=5138001 file=extracted/1008/5138001/제26회_감정평가사 1차_기출/제26회_감정평가사 1차_기출/1교시_감평_A형.pdf
%% conversion: kordoc 4.24.51 -> tools/gen1.py (auto); review=proofread-sonnet; flags=none
\begin{question}{70}
한 나라의 거시경제모형이 다음과 같을 때, 생산물시장과 화폐시장이 동시에 균형을 이루는 국민소득은? (단, $Y$는 국민소득, $r$은 이자율, $Y_{D}$는 가처분소득이다.)
\begin{framedtext}
소비함수: $C=150+0.8Y_{D}$ \par
투자함수: $I=100-100r$ \par
정부지출: $G=250$ \par
조\quad 세: $T=0.25Y$ \par
화폐수요: $M_{D}=50+0.2Y-200r$ \par
화폐공급: $M_{S}=250$ \par
\end{framedtext}
\begin{choices}
  \item 250
  \item 1,200
  \item 1,250
  \item 1,500
  \item 3,000
\end{choices}
\end{question}
\answer{2}
